\chapter{Literature Review}

This chapter examines the protocols, products and enabling technologies that
define the current state of remote desktop access, and closes by identifying the
gap that \updesk{} occupies.

\section{Remote Framebuffer and VNC}

Virtual Network Computing, developed at the Olivetti Research Laboratory and
standardised as the Remote Framebuffer protocol in RFC~6143, is the conceptual
ancestor of most open remote desktop software. Its model is deliberately thin:
the server exports a rectangular framebuffer, the client requests updates, and
the server responds with rectangles that have changed, optionally encoded using
schemes such as Hextile, ZRLE or Tight.

The strength of this design is portability. Because the protocol reasons about
pixels rather than drawing commands, a VNC server can be written for any system
that can read its own screen. The corresponding weakness is efficiency. Video
playback or window dragging invalidates large regions, and rectangle-oriented
encodings degrade badly under such motion. RFB also assumes the client can open
a TCP connection to the server, which on the modern Internet is usually false
without manual port forwarding. Security was not part of the original design;
encryption is provided, when at all, by tunnelling or by vendor extensions.

\section{Microsoft Remote Desktop Protocol}

RDP takes the opposite approach. Rather than shipping pixels, it intercepts
drawing operations and transmits high-level primitives --- glyphs, blits,
geometric fills --- which the client replays locally. For conventional desktop
work this is dramatically more efficient than any framebuffer encoding, and
successive revisions have added RemoteFX and H.264 codecs for content that does
not decompose into primitives.

RDP is mature, well integrated with Windows session management, and capable of
reaching the login screen because it is part of the operating system rather than
an application running inside a session. Its limitations for the present purpose
are that it is tightly coupled to the Windows display architecture, that its
implementation is closed, and that reaching an RDP host from outside a network
requires either a routable address, a virtual private network, or a Remote
Desktop Gateway. Exposing RDP directly to the Internet is widely documented as a
significant risk and is a common initial vector for ransomware intrusion.

\section{Commercial Peer-Assisted Products}

TeamViewer, AnyDesk and similar commercial tools solved the reachability problem
in the way users actually experience it: the host registers with vendor
infrastructure and is addressed by a short numeric identifier, and the vendor
arranges connectivity irrespective of NAT. AnyDesk's DeskRT codec, in
particular, is specialised for graphical user interface content and achieves
notably low latency.

These products are highly capable, and their user experience informed the design
of \updesk{}'s connect-identifier scheme. The architectural objection is one of
trust rather than quality. Session traffic is arranged, and frequently relayed,
by infrastructure the user neither controls nor can audit, and the identifier
namespace itself is vendor-operated. Published incidents in which support
sessions were abused illustrate that the concentration of capability is a real
consideration and not merely a theoretical one.

\section{RustDesk and the Open Alternative}

RustDesk is an open-source remote desktop implementation written in Rust that
mirrors the commercial user experience --- numeric identifiers, self-hostable
relay and rendezvous servers --- while permitting inspection of the source. Its
architecture separates a native capture library from a service component, which
is precisely the structure required to reach privileged desktops on Windows.

RustDesk is the closest prior work to this project and was studied carefully
during design. \updesk{} differs in two respects that are relevant to the
contribution claimed here. First, it uses standard WebRTC as its transport, so
media encryption, congestion control and ICE behaviour are inherited from a
widely reviewed implementation rather than from a bespoke protocol. Second, its
signaling layer authenticates every endpoint with an Ed25519
challenge--response, with enrollment gated by one-time codes, making device
identity a first-class element of the protocol rather than a property of a
shared password.

\section{WebRTC as a Transport}

WebRTC is a set of standards for real-time communication between peers,
comprising media capture, a peer connection abstraction, and data channels. Its
relevance here rests on three properties.

\textbf{Connectivity.} Interactive Connectivity Establishment (RFC~8445) has
each peer gather candidate transport addresses --- host addresses, server-
reflexive addresses discovered through STUN, and relayed addresses obtained from
TURN --- exchange them, and probe candidate pairs systematically until a
working path is found. This resolves the NAT traversal problem that defeats VNC
and RDP without operator configuration.

\textbf{Security.} Media encryption is not optional. Peers perform a DTLS
handshake and derive SRTP keys from it, so the payload is protected end to end.
A signaling server that only relays session descriptions and candidates cannot
recover those keys, which is exactly the property \updesk{} requires.

\textbf{Adaptivity.} The stack includes congestion control, packet loss
concealment and jitter buffering developed for conferencing traffic, which
transfers directly to screen streaming.

WebRTC deliberately does not specify signaling. Peers must exchange SDP offers,
answers and ICE candidates over some out-of-band channel of the application's
choosing, and the design of that channel is where an application's security
model is largely decided.

\section{Screen Capture on Windows}

Three mechanisms are relevant. The \texttt{getDisplayMedia} API available to web
content is simple and cross-platform, but the browser security model requires a
user gesture and a source-selection dialogue, and content on the secure desktop
is never visible to it. The Desktop Duplication API exposed through DXGI
provides efficient access to the composited desktop surface with damage regions,
and is the basis of most native Windows capture. The newer
Windows.Graphics.Capture API offers per-window capture with a cleaner interface
at the cost of requiring recent Windows versions.

Critically, none of these can observe the Winlogon secure desktop from an
ordinary user session. Windows isolates the login screen and UAC prompts on a
separate desktop in a distinct session precisely to prevent one application from
observing or driving another's credential entry. Reaching it requires a process
running as SYSTEM that can attach to the active desktop --- which is why the
architecture in Chapter~4 includes a service component.

\section{Identified Gap}

The survey suggests the following. Open protocols such as RFB and RDP are
transparent but assume reachability that ordinary networks do not provide.
Commercial products provide reachability but require trust in infrastructure
that cannot be audited. RustDesk narrows the gap substantially but implements
its own transport and its own security layer. There remains room for a system
that (i) uses standard, widely reviewed WebRTC for transport and encryption,
(ii) authenticates endpoints with public-key identity so that the signaling
server can verify who is connecting without being able to read what is sent, and
(iii) treats unattended, pre-login, secure-desktop-capable operation as a design
requirement rather than an afterthought. \updesk{} is an attempt to occupy that
position.

\tabref{tab:comparison} summarises the comparison.

\begin{table}[H]
\centering
\caption{Comparison of remote desktop approaches}
\label{tab:comparison}
\small
\begin{tabular}{@{}lccccc@{}}
\toprule
\textbf{System} & \textbf{Open} & \textbf{NAT} & \textbf{E2E} & \textbf{Pre-login} & \textbf{Transport} \\
 & & \textbf{traversal} & \textbf{encrypted} & \textbf{access} & \\
\midrule
VNC (RFB)   & Yes & No      & Optional & No  & TCP \\
RDP         & No  & Gateway & Yes      & Yes & TCP/UDP \\
TeamViewer  & No  & Yes     & Claimed  & Yes & Proprietary \\
AnyDesk     & No  & Yes     & Claimed  & Yes & Proprietary \\
RustDesk    & Yes & Yes     & Yes      & Yes & Custom \\
\textbf{UpDesk} & Yes & Yes & Yes      & Yes & WebRTC \\
\bottomrule
\end{tabular}
\end{table}
